\section{Limitaciones del modelo y justificación del enfoque}

\subsection{Limitaciones metodológicas}

Todo modelo de forecasting tiene limitaciones. En el software ICOCIV, las principales son:

\begin{itemize}
    \item \textbf{Dependencia de datos históricos}: el modelo aprende de patrones pasados. Si el futuro cambia estructuralmente, la proyección puede perder precisión.
    \item \textbf{Horizonte de predicción}: la confiabilidad disminuye cuando el horizonte se aleja del último dato observado.
    \item \textbf{Variables omitidas}: el modelo no incorpora explícitamente tasa de cambio, precios internacionales, política monetaria ni variables macroeconómicas.
    \item \textbf{Choques externos}: eventos excepcionales pueden alterar costos de forma no anticipada.
    \item \textbf{Cambios metodológicos}: cualquier modificación en la operación estadística o clasificación puede afectar comparabilidad.
    \item \textbf{Autocorrelación residual}: si persiste en los errores, indica que la regresión no capturó toda la dinámica temporal.
\end{itemize}

Reconocer estas limitaciones fortalece el documento. Un sistema técnicamente serio no oculta la incertidumbre; la comunica para que el usuario tome decisiones informadas.

\subsection{Por qué usar modelos interpretables}

El software prioriza modelos interpretables porque su propósito es institucional y académico. En ingeniería civil, las proyecciones pueden respaldar informes, discusiones contractuales o análisis de variación de costos. En esos contextos, es necesario explicar el razonamiento del resultado.

Los modelos de regresión permiten responder preguntas como:

\begin{itemize}
    \item ¿La tendencia es aproximadamente lineal?
    \item ¿Existe curvatura o aceleración?
    \item ¿El modelo ajusta bien sin complejidad excesiva?
    \item ¿Los residuos muestran autocorrelación?
    \item ¿Cómo fue el error en backtesting?
\end{itemize}

Estas respuestas son comprensibles para un lector con formación estadística básica o intermedia. Además, pueden documentarse en reportes técnicos con ecuaciones, métricas y trazabilidad.

\subsection{Por qué no usar inicialmente modelos de caja negra}

Modelos como redes neuronales, ensambles complejos o algoritmos de machine learning de alta capacidad pueden ser útiles en ciertos problemas predictivos. Sin embargo, no son la mejor primera opción para este proyecto por varias razones:

\begin{itemize}
    \item las series ICOCIV disponibles son relativamente cortas para modelos de alta complejidad;
    \item el objetivo exige explicación y trazabilidad, no solo precisión;
    \item los usuarios técnicos deben comprender la lógica de selección;
    \item el software debe funcionar localmente en computadores de bajos recursos;
    \item la validación académica requiere justificar supuestos y limitaciones.
\end{itemize}

La decisión no es un rechazo general a la inteligencia artificial, sino una elección metodológica proporcional al problema. Para un trabajo de grado aplicado a índices económicos de construcción, la interpretabilidad y la reproducibilidad son criterios centrales.

\subsection{Relación entre implementación y metodología}

\begin{longtable}{p{0.25\textwidth}p{0.32\textwidth}p{0.33\textwidth}}
\caption{Relación entre decisiones del software y justificación metodológica}
\label{tab:implementacion_metodologia}\\
\toprule
\textbf{Componente} & \textbf{Implementación} & \textbf{Justificación} \\
\midrule
\endfirsthead
\toprule
\textbf{Componente} & \textbf{Implementación} & \textbf{Justificación} \\
\midrule
\endhead
Lectura de periodos & Detección dinámica desde encabezados reales & Evita series artificiales y errores de calendario. \\
Validación de calidad & Continuidad, duplicados, faltantes, longitud mínima & Asegura base temporal confiable antes de modelar. \\
Modelos & Lineal, polinómico, logarítmico y robusto & Mantienen interpretabilidad y flexibilidad moderada. \\
Selección & AICc, residuos, \(R^2\), backtesting y coherencia económica & Evita depender de una sola métrica. \\
Backtesting & Origen rodante con orden temporal & Simula el uso real pasado-futuro. \\
Reportes & DOCX con métricas, gráfica y trazabilidad & Convierte el resultado en evidencia documental auditable. \\
\bottomrule
\end{longtable}

\subsection{Escalabilidad metodológica}

El enfoque actual permite extensiones futuras sin romper la lógica académica. Entre ellas se encuentran:

\begin{itemize}
    \item incorporación de variables macroeconómicas en regresiones explicativas;
    \item análisis de cambios estructurales;
    \item intervalos de predicción con métodos más formales;
    \item ventanas rodantes para series con cambios de régimen.
\end{itemize}

Estas extensiones deben incorporarse de manera gradual y documentada. La prioridad sigue siendo que cada mejora preserve trazabilidad, validación temporal y claridad para el usuario técnico.
